\section{Finite separation and rank transfer}\label{sec:separation}

\begin{lemma}[One finite separator]\label{lem:separator}
\notformalized
Under \eqref{eq:separation}, there is a finite quotient
$q:\Gamma\twoheadrightarrow Q$ in which every matched port pair is
still disconnected. This quotient need not be a $p$-group or a member
of the chain \eqref{eq:chain}.
\end{lemma}
\begin{informalproof}
The finite base incidence multigraph has ports $L_v,R_v$ and one
edge per triple. Orient a traversal from $L_u$ to $R_v$ with label
$g$ and its reverse with label $g^{-1}$. The voltage of a path with
chronological labels $a_1,\ldots,a_m$ is $a_m\cdots a_1$.

In a base component $D$, choose a root and a spanning tree. Let $t_z$
be the voltage of the root-to-$z$ tree path and let $J_D$ be the
subgroup of voltages of closed root walks. It is generated by the
finitely many elements $t_y^{-1}a t_x$ from oriented edges
$x\xrightarrow{a}y$: writing $a=t_y(t_y^{-1}a t_x)t_x^{-1}$
makes the tree labels telescope along a closed walk. The path
voltages from $x$ to $y$ are exactly $t_yJ_Dt_x^{-1}$, by closing
a path at the root, and conversely inserting a closed root walk.

If $L_v,R_v$ are in different base components, their lifts are always
disconnected. Otherwise their matched lifts in a quotient are connected
exactly when
\[
 q(t_{R_v}^{-1}t_{L_v})\in q(J_D).
\]
Thus \eqref{eq:separation} gives a finite list of pairs $(J,g)$ with
$J$ finitely generated and $g\notin J$.

Here is the required finite separation argument. In the coset graph
with vertices $zJ$, retain the paths from $J$ spelled by a finite
generating set of $J$ and by $g$. For a left action, traverse a word
from its rightmost letter. Their finite vertex union $S$ contains
both $J$ and the distinct vertex $gJ$. Each free generator acts on
the retained edges by a partial injection of $S$ into itself. Complete
it to a permutation, matching the unused domain and range vertices.
Freeness extends these permutations to a homomorphism
$\Gamma\to\operatorname{Sym}(S)$. The retained generator paths of
$J$ fix the vertex $J$, while $g$ does not. The image of $g$ is
therefore outside the image of $J$.

Take the product of these finitely many separating homomorphisms and
let $Q$ be its image. Each forbidden membership fails in its designated
factor, hence in $Q$. If the list is empty, use the trivial quotient.
\end{informalproof}

This is the standard subgroup-separation mechanism specialized to the
two ports of every type. Its role is to supply a finite rank witness;
it does not replace the actual chain by this auxiliary quotient.

\begin{proposition}[A finite rank witness]\label{prop:inner-rank}
\notformalized
For any field $K$ and widths of total $B>0$, let
$M=I-\sum_eT_eg_e\in\operatorname{Mat}_B(K[\Gamma])$ fit a template
satisfying \eqref{eq:separation}. Then
\[
 \irk_{K[\Gamma]}M\ge\lceil B/2\rceil.
\]
Here inner rank is the minimum $k$ for a factorization through $k$
columns over the group algebra.
\end{proposition}
\begin{informalproof}
Use the quotient $Q$ from \cref{lem:separator}. Order its incidence
components uniformly at random. Select a coordinate $(v,x)$ when
its left component precedes its right component, and give it weight
$w_v$. The components differ, so its selection probability is $1/2$.
Some order selects total weight at least $B|Q|/2$.

For a selected input $\alpha$ and selected receiver $\beta$ joined
by a dependency, the left component of $\alpha$ equals the right
component of $\beta$. Selection of $\beta$ therefore places its left
component strictly before that of $\alpha$. The selected principal
block of $M_Q$, ordered by left component, is triangular with identity
diagonal. Consequently $\rank_K M_Q\ge B|Q|/2$.

A factorization $M=XY$ through $k$ columns specializes under the
regular representation to a factorization of $M_Q$ through $k|Q|$
columns. Hence $k|Q|\ge B|Q|/2$ and $k\ge\lceil B/2\rceil$.
No semisimplicity is needed, even when $p$ divides $|Q|$.
\end{informalproof}

\paragraph{Published algebraic inputs.}
Cohn's theorem identifies inner rank over a free-group algebra with
rank in its universal division ring~\cite{Cohn,JaikinSouza}.
For the subgroup and chain in \eqref{eq:chain},
Jaikin-Zapirain--Souza~\cite[Theorem~8.1]{JaikinSouza} gives, for
each fixed matrix $M$ over $\F_p[\Gamma]$,
\begin{equation}\label{eq:rank-approx}
 \lim_i\frac{\rank_{\F_p}M_{Q_i}}{q_i}
       =\irk_{\F_p[\Gamma]}M.
\end{equation}
These are imported theorems, not formalized here. The source uses a
right action on rows. Its finite rank equals our left-action rank on
columns: the pairing $\langle y,x\rangle=\tau(\sum_jy_jx_j)$,
where $\tau$ extracts the identity coefficient in $\F_p[Q_i]$,
is nondegenerate and satisfies $\langle y,Mx\rangle=\langle yM,x\rangle$.

\begin{theorem}[Fixed prime-field checks]\label{thm:prime-field}
\notformalized
With \eqref{eq:chain} and \eqref{eq:separation}, every fixed
$\F_p$ check of total width $B>0$ satisfies
\[
 \lim_i\frac{\dim_{\F_p}\ker M_{Q_i}}{Bq_i}
       \le\frac{\lfloor B/2\rfloor}{B}.
\]
The same upper bound on the limsup holds when the $\F_p$ coefficients
vary with $i$, while the template and widths remain fixed.
\end{theorem}
\begin{informalproof}
Combine \cref{prop:inner-rank} with \eqref{eq:rank-approx} and
rank-nullity. There are finitely many $\F_p$ assignments to the fixed
coefficient entries, so the maximum of their convergence thresholds
works for every assignment at once.
\end{informalproof}

For a fixed finite extension $K/\F_p$ of degree $a$, restriction of
scalars multiplies both kernel dimension and $B$ by $a$. Thus the
pointwise half bound holds over $K$, with the possibly stronger bound
$\lfloor aB/2\rfloor/(aB)$. This argument alone does not yield
$\lfloor B/2\rfloor/B$ over $K$ or a uniform threshold over all $a$.
The next section supplies the latter uniformity at the half constant.

\begin{lemma}[Long paired distances imply infinite separation]
\label{lem:distance-separation}
\notformalized
For a fixed template with $n=|V|>0$, suppose finite quotient lifts
satisfy, for every fixed integer $r\ge0$,
\[
 \#\{(v,x):\dist_{I_{Q_i}}(L_{(v,x)},R_{(v,x)})\le r\}
       =o(nq_i).
\]
Then \eqref{eq:separation} holds. This implication needs no nesting
or prime-power condition on the quotients.
\end{lemma}
\begin{informalproof}
An infinite-cover matched path at one type has finite length $r$.
It descends to a walk of that length in every quotient. Right
translation supplies such a walk at all $q_i$ sheets of that type,
giving proportion at least $1/n$, a contradiction.
\end{informalproof}

Thus the native paired-distance invariant can supply the separation
premise \emph{if} a fixed-template cover representation is already
available. The rank proof provides no such representation.
